\section{Design do Experimento}

\begin{table}[h]
\centering
\small
\begin{tabular}{|l|c|c|}
\hline
\textbf{Aspecto} & \textbf{Grupo Controle} & \textbf{Grupo Tratamento} \\
\hline
\textit{Estratégia de Seleção} & Aleatória ou alfabética (atual) & PhoneScorer \\
\hline
\textit{Critério de Alocação} & \multicolumn{2}{c|}{Aleatório por CPF (50/50)} \\
\hline
\textit{Tamanho} & \multicolumn{2}{c|}{1.335 CPFs por grupo} \\
\hline
\textit{Período} & \multicolumn{2}{c|}{7 dias consecutivos} \\
\hline
\end{tabular}
\caption{Estrutura do Teste A/B: alocação aleatória de CPFs a controle e tratamento.}
\label{tab:design_ab}
\end{table}

\subsection*{Definições do Experimento}

\begin{itemize}
    \item \textbf{Unidade de Randomização:} CPF (não telefone). Cada CPF é alocado uma única vez a um dos grupos no início do experimento.

    \item \textbf{Persistência da Alocação:} A alocação é fixa ao longo de todo o experimento (\textit{sticky assignment}), garantindo que um mesmo CPF não transite entre controle e tratamento.

    \item \textbf{Unidade de Análise:} As métricas são agregadas no nível de CPF.

    \item \textbf{Regra de Exposição:} Para cada CPF, o grupo tratamento (PhoneScorer) seleciona o \textbf{telefone com maior score} (top-1) para envio da mensagem. O grupo controle segue a estratégia atual (aleatória ou alfabética).

    \item \textbf{Janela de Observação:} Para cada mensagem enviada, o status de entrega e leitura é avaliado após uma janela fixa de 24 horas.

    \item \textbf{Critérios de Inclusão:} A análise considera apenas mensagens com janela de observação completa de 24 horas. CPFs sem envios no período não são incluídos no cálculo das métricas.

    \item \textbf{Balanceamento:} A alocação é realizada de forma balanceada entre os grupos (50/50), resultando em aproximadamente 1.335 CPFs por grupo.

    \item \textbf{Independência:} Cada CPF é tratado como uma observação independente no experimento.
\end{itemize}